% LuaLaTeX + Metropolis Beamer — 그림 적극 활용 버전
\documentclass[aspectratio=169,11pt]{beamer}

\usetheme[progressbar=frametitle]{metropolis}

% === 패키지 ===
\usepackage{fontspec}
\usepackage[hangul]{luatexko}
\usepackage{graphicx}
\usepackage{xcolor}
\usepackage{booktabs}
\usepackage{array}
\usepackage{emoji}
\usepackage{tikz}
\usepackage{makecell}

% === 폰트 ===
\setmainfont[Ligatures=TeX]{TeX Gyre Heros}
\setsansfont{TeX Gyre Heros}
\setmonofont{Noto Sans Mono}
\setmainhangulfont{Noto Sans CJK KR}[BoldFont={Noto Sans CJK KR Bold}]
\setsanshangulfont{Noto Sans CJK KR}[BoldFont={Noto Sans CJK KR Bold}]
\setemojifont{Noto Color Emoji}

% === 색상 ===
\definecolor{mainblue}{HTML}{1976D2}
\definecolor{mainteal}{HTML}{00897B}
\definecolor{mainorange}{HTML}{F57C00}
\definecolor{mainpurple}{HTML}{7B1FA2}
\definecolor{maingreen}{HTML}{388E3C}
\definecolor{lightblue}{HTML}{E3F2FD}
\definecolor{lightmint}{HTML}{E0F2F1}
\definecolor{lightcream}{HTML}{FFF8E1}

\setbeamercolor{frametitle}{bg=mainblue, fg=white}
\setbeamercolor{progress bar}{fg=mainorange}

% === 표 열 타입 ===
\newcolumntype{C}[1]{>{\centering\arraybackslash}m{#1}}
\newcolumntype{L}[1]{>{\raggedright\arraybackslash}m{#1}}

% === 이미지 경로 ===
\graphicspath{{images/}}

% === PDF 설정 ===
\hypersetup{unicode=true, pdfencoding=unicode}

\title{AI디지털리터러시}
\subtitle{1주차 2차시 — 기계학습의 원리}
\author{청주교육대학교}
\date{2026년 2학기}

\begin{document}

% ============================================================
% 슬라이드 1: 타이틀
% ============================================================
\maketitle

% ============================================================
% 슬라이드 2: 목차
% ============================================================
\begin{frame}{오늘의 수업 흐름}
\begin{description}[labelwidth=3cm]
\item[\emoji{clapper-board} \textbf{도입} (5분)] AI, ML, DL의 포함 관계
\item[\emoji{blue-book} \textbf{전개 1} (15분)] 기계학습의 정의
\item[\emoji{green-book} \textbf{전개 2} (15분)] 3가지 학습 유형 (지도/비지도/강화)
\item[\emoji{orange-book} \textbf{전개 3} (10분)] 머신러닝 4단계 프로세스
\item[\emoji{memo} \textbf{정리} (5분)] 엔트리 AI와의 연결
\end{description}
\end{frame}

% ============================================================
% 슬라이드 3: 복습
% ============================================================
\begin{frame}{지난 시간 핵심 내용}
\emoji{check-mark-button} \textbf{AI의 정의} — 인간의 지능을 모방하는 컴퓨터 시스템

\vspace{0.4cm}
\emoji{check-mark-button} \textbf{약한 AI vs 강한 AI} — 현재 모든 AI는 약한 AI (특정 작업 전문가)

\vspace{0.4cm}
\emoji{check-mark-button} \textbf{AI의 4가지 기능} — 인식, 예측, 분류, 생성

\vspace{0.4cm}
\emoji{check-mark-button} \textbf{AI 발전사} — AI의 봄과 겨울 반복 $\rightarrow$ 현재 제3의 봄
\end{frame}

% ============================================================
% 슬라이드 4: 핵심 질문
% ============================================================
\begin{frame}{오늘의 핵심 질문}
\begin{center}
\vspace{0.5cm}
{\Large \emoji{light-bulb} ``1차시에서 배운 AI,}

\vspace{0.3cm}
{\Large 그 안에는 무엇이 있을까요?''}

\vspace{1cm}
\textbf{오늘의 목표}

\vspace{0.3cm}
\emoji{check-mark-button} 기계학습의 원리 이해하기

\emoji{check-mark-button} 3가지 학습 유형 구분하기
\end{center}
\end{frame}

% %%%%%%%%%%%%%%%%%%%%%%%%%%%%%%%%%%%%%%%%%%%%%%%%%%%%%%%%%%%%
% 도입: AI, ML, DL의 포함 관계
% %%%%%%%%%%%%%%%%%%%%%%%%%%%%%%%%%%%%%%%%%%%%%%%%%%%%%%%%%%%%
\section{도입: AI, ML, DL의 포함 관계}

% ============================================================
% 슬라이드 5: 그림 01 — AI⊃ML⊃DL (그림 단독)
% ============================================================
\begin{frame}{AI, ML, DL의 포함 관계}
\begin{center}
\includegraphics[width=0.85\textwidth]{AI디지털리터러시_1주차1차시_06_AI_ML_DL포함관계_v1.0.1.0_점선삭제_2026-02-08_0010.png}

\vspace{0.2cm}
{\tiny 그림 01: AI\_ML\_DL포함관계}
\end{center}
\end{frame}

% ============================================================
% 슬라이드 6: AI, ML, DL 핵심 설명
% ============================================================
\begin{frame}{AI, ML, DL — 핵심 정리}
\textbf{AI (Artificial Intelligence)} — 인간의 지능을 모방하는 컴퓨터 시스템
\begin{itemize}
\item 가장 넓은 개념
\end{itemize}

\vspace{0.3cm}
\textbf{ML (Machine Learning)} — 데이터로부터 패턴을 학습하는 기술
\begin{itemize}
\item AI를 구현하는 대표적 방법
\end{itemize}

\vspace{0.3cm}
\textbf{DL (Deep Learning)} — 인공신경망을 사용하는 기계학습
\begin{itemize}
\item ML의 한 종류 (최근 주목)
\end{itemize}

\vspace{0.3cm}
\emoji{point-right} \textbf{오늘 초점}: 기계학습(ML)의 원리와 3가지 학습 유형
\end{frame}

% %%%%%%%%%%%%%%%%%%%%%%%%%%%%%%%%%%%%%%%%%%%%%%%%%%%%%%%%%%%%
% 전개 1: 기계학습의 정의
% %%%%%%%%%%%%%%%%%%%%%%%%%%%%%%%%%%%%%%%%%%%%%%%%%%%%%%%%%%%%
\section{전개 1: 기계학습의 정의}

% ============================================================
% 슬라이드 7: 기계학습이란?
% ============================================================
\begin{frame}{기계학습이란?}
\begin{center}
\vspace{0.3cm}
{\large \textbf{Machine Learning (기계학습)}}

\vspace{0.5cm}
\begin{tikzpicture}
\node[draw=mainblue, fill=lightblue, rounded corners=8pt, inner sep=12pt, text width=10cm, align=center] {
{\large \emoji{light-bulb} 데이터로부터 패턴을 학습하여}\\[0.2cm]
{\large 새로운 상황에서 예측하는 AI 기술}
};
\end{tikzpicture}

\vspace{0.5cm}
\emoji{bar-chart} \textbf{데이터} \hspace{1cm}
\emoji{magnifying-glass-tilted-right} \textbf{패턴 발견} \hspace{1cm}
\emoji{bullseye} \textbf{예측}
\end{center}
\end{frame}

% ============================================================
% 슬라이드 8: 전통적 프로그래밍 vs 기계학습
% ============================================================
\begin{frame}{전통적 프로그래밍 vs 기계학습}
\begin{columns}
\begin{column}{0.45\textwidth}
\begin{center}
\textbf{\color{mainorange}전통적 프로그래밍}

\vspace{0.3cm}
\begin{tikzpicture}
\node[draw=mainorange, fill=lightcream, rounded corners=5pt, inner sep=8pt, text width=4.5cm, align=center] {
규칙 + 데이터 $\rightarrow$ 결과\\[0.2cm]
사람이 규칙 작성 \emoji{writing-hand}\\
컴퓨터가 실행 \emoji{desktop-computer}
};
\end{tikzpicture}
\end{center}
\end{column}
\begin{column}{0.45\textwidth}
\begin{center}
\textbf{\color{mainblue}기계학습}

\vspace{0.3cm}
\begin{tikzpicture}
\node[draw=mainblue, fill=lightblue, rounded corners=5pt, inner sep=8pt, text width=4.5cm, align=center] {
데이터 + 결과 $\rightarrow$ 규칙 발견\\[0.2cm]
사람이 데이터 제공 \emoji{bar-chart}\\
컴퓨터가 규칙 발견 \emoji{robot}
};
\end{tikzpicture}
\end{center}
\end{column}
\end{columns}
\end{frame}

% ============================================================
% 슬라이드 9: 왜 기계학습이 필요한가?
% ============================================================
\begin{frame}{왜 기계학습이 필요한가?}
\begin{columns}
\begin{column}{0.45\textwidth}
\begin{center}
\textbf{\color{mainorange}\emoji{cross-mark} 규칙 작성이 불가능}

\vspace{0.3cm}
고양이 vs 개 구분 규칙?

\vspace{0.2cm}
귀 모양? 털 색깔? 크기?

\vspace{0.2cm}
$\rightarrow$ 예외가 너무 많음!
\end{center}
\end{column}
\begin{column}{0.45\textwidth}
\begin{center}
\textbf{\color{maingreen}\emoji{check-mark-button} 기계학습의 해결}

\vspace{0.3cm}
수천 장의 사진 제공

\vspace{0.2cm}
컴퓨터가 스스로 패턴 발견

\vspace{0.2cm}
$\rightarrow$ 새로운 사진도 구분!
\end{center}
\end{column}
\end{columns}
\end{frame}

% ============================================================
% 슬라이드 10: 일상 속 기계학습
% ============================================================
\begin{frame}{일상 속 기계학습}
\emoji{e-mail} \textbf{스팸 필터} — 스팸 메일 패턴 학습 $\rightarrow$ 자동 분류

\vspace{0.4cm}
\emoji{clapper-board} \textbf{영화 추천} — 내가 본 영화 분석 $\rightarrow$ 비슷한 영화 추천

\vspace{0.4cm}
\emoji{microphone} \textbf{음성 인식} — 다양한 목소리 학습 $\rightarrow$ 텍스트 변환

\vspace{0.4cm}
\emoji{grinning-face} \textbf{얼굴 인식} — 수많은 얼굴 학습 $\rightarrow$ 특정 인물 인식
\end{frame}

% ============================================================
% 슬라이드 11: 기계학습의 핵심 3요소
% ============================================================
\begin{frame}{기계학습의 핵심 3요소}
\begin{center}
\vspace{0.3cm}
\begin{tikzpicture}
\node[draw=mainblue, fill=lightblue, rounded corners=8pt, inner sep=15pt, minimum width=3cm, minimum height=1.5cm, text width=2.8cm, align=center] at (0,0) {\textbf{① 데이터}\\학습의 재료};
\node at (2.5,0) {\Large $\Rightarrow$};
\node[draw=mainteal, fill=lightmint, rounded corners=8pt, inner sep=15pt, minimum width=3cm, minimum height=1.5cm, text width=2.8cm, align=center] at (5,0) {\textbf{② 학습}\\패턴을 찾는 과정};
\node at (7.5,0) {\Large $\Rightarrow$};
\node[draw=mainpurple, fill=lightcream, rounded corners=8pt, inner sep=15pt, minimum width=3cm, minimum height=1.5cm, text width=2.8cm, align=center] at (10,0) {\textbf{③ 예측}\\새로운 상황에 적용};
\end{tikzpicture}
\end{center}
\end{frame}

% %%%%%%%%%%%%%%%%%%%%%%%%%%%%%%%%%%%%%%%%%%%%%%%%%%%%%%%%%%%%
% 전개 2: 3가지 학습 유형
% %%%%%%%%%%%%%%%%%%%%%%%%%%%%%%%%%%%%%%%%%%%%%%%%%%%%%%%%%%%%
\section{전개 2: 3가지 학습 유형}

% ============================================================
% 슬라이드 12: 그림 02 — 3가지 학습 유형 비교 (그림 단독)
% ============================================================
\begin{frame}{3가지 학습 유형 비교}
\begin{center}
\includegraphics[width=0.85\textwidth]{AI디지털리터러시_1주차2차시_10_3가지학습유형비교.png}

\vspace{0.1cm}
{\tiny 그림 02: 3가지\_학습유형비교}
\end{center}
\end{frame}

% ============================================================
% 슬라이드 13: 지도학습 정의
% ============================================================
\begin{frame}{① 지도학습 (Supervised Learning)}
\begin{center}
\begin{tikzpicture}
\node[draw=mainblue, fill=lightblue, rounded corners=8pt, inner sep=12pt, text width=10cm, align=center] {
\emoji{light-bulb} \textbf{정답(라벨)이 있는} 데이터로 학습
};
\end{tikzpicture}
\end{center}

\vspace{0.3cm}
\textbf{비유}: 선생님이 정답을 알려주며 가르치는 것

\vspace{0.3cm}
\textbf{데이터 구조}:
\begin{itemize}
\item 고양이 사진 $\rightarrow$ ``고양이'' (정답)
\item 개 사진 $\rightarrow$ ``개'' (정답)
\end{itemize}

\vspace{0.3cm}
\emoji{point-right} 엔트리 이미지 분류 모델이 바로 \textbf{지도학습}!
\end{frame}

% ============================================================
% 슬라이드 14: 분류 vs 회귀 비교
% ============================================================
\begin{frame}{지도학습의 2가지 유형: 분류 vs 회귀}
\begin{center}
\begin{tabular}{C{1.8cm} L{4cm} L{4cm}}
\toprule
\textbf{구분} & \textbf{회귀 (얼마나?)} & \textbf{분류 (어느 것?)} \\
\midrule
\textbf{답 형태} & 숫자 (연속적) & 범주 (구분) \\
\textbf{예시 1} & 시험 점수: 85점 & 합격/불합격: 합격 \\
\textbf{예시 2} & 집값: 3억 5천 & 아파트/빌라: 아파트 \\
\textbf{예시 3} & 기온: 18도 & 더위/추위: 더위 \\
\bottomrule
\end{tabular}
\end{center}
\end{frame}

% ============================================================
% 슬라이드 15: 그림 03 — 지도 vs 비지도 데이터 구조 (그림 단독)
% ============================================================
\begin{frame}{② 비지도학습 — 지도학습과 무엇이 다를까?}
\begin{center}
\includegraphics[width=0.85\textwidth]{AI디지털리터러시_1주차2차시_15_지도학습vs비지도학습_데이터구조비교.png}

\vspace{0.1cm}
{\tiny 그림 03: 지도학습vs비지도학습\_데이터구조비교}
\end{center}
\end{frame}

% ============================================================
% 슬라이드 16: 그림 04 — 라벨 없는 3가지 이유 (그림 단독)
% ============================================================
\begin{frame}{왜 라벨(정답) 없이 학습할까?}
\begin{center}
\includegraphics[width=0.85\textwidth]{AI디지털리터러시_1주차2차시_17_라벨없는_3가지이유.png}

\vspace{0.1cm}
{\tiny 그림 04: 라벨없는\_3가지이유}
\end{center}
\end{frame}

% ============================================================
% 슬라이드 17: 그림 05 — 군집화 3단계 과정 (그림 단독)
% ============================================================
\begin{frame}{비지도학습의 핵심: 군집화 (Clustering)}
\begin{center}
\includegraphics[width=0.85\textwidth]{AI디지털리터러시_1주차2차시_16_군집화_3단계과정.png}

\vspace{0.1cm}
{\tiny 그림 05: 군집화\_3단계과정}
\end{center}
\end{frame}

% ============================================================
% 슬라이드 18: 그림 06 — 고객 세분화 결과 (columns: 그림+키워드)
% ============================================================
\begin{frame}{비지도학습 예시: 고객 세분화}
\begin{columns}
\begin{column}{0.55\textwidth}
\includegraphics[width=\textwidth]{AI디지털리터러시_1주차2차시_18_고객세분화결과_4개그룹.png}

{\tiny 그림 06: 고객세분화결과\_4개그룹}
\end{column}
\begin{column}{0.40\textwidth}
\textbf{쇼핑몰 고객 10,000명}

\vspace{0.3cm}
AI가 4개 그룹 자동 분류:
\begin{itemize}
\item \emoji{money-bag} 알뜰형
\item \emoji{gem-stone} 명품형
\item \emoji{fire} 트렌드형
\item \emoji{package} 실속형
\end{itemize}
\end{column}
\end{columns}
\end{frame}

% ============================================================
% 슬라이드 19: 그림 07 — 강화학습 루프 (그림 단독)
% ============================================================
\begin{frame}{③ 강화학습 — 시행착오로 배우는 AI}
\begin{center}
\includegraphics[width=0.85\textwidth]{AI디지털리터러시_1주차2차시_14_강화학습루프.png}

\vspace{0.1cm}
{\tiny 그림 07: 강화학습루프 (Agent $\leftrightarrow$ Environment)}
\end{center}
\end{frame}

% ============================================================
% 슬라이드 20: 강화학습 핵심 설명
% ============================================================
\begin{frame}{강화학습 (Reinforcement Learning)}
\begin{center}
\begin{tikzpicture}
\node[draw=mainpurple, fill=lightcream, rounded corners=8pt, inner sep=12pt, text width=10cm, align=center] {
\emoji{light-bulb} \textbf{시행착오}를 통해 \textbf{보상을 최대화}하는 학습
};
\end{tikzpicture}
\end{center}

\vspace{0.4cm}
\textbf{핵심 개념}:
\begin{itemize}
\item \textbf{에이전트}(Agent): AI — 행동을 결정하는 주체
\item \textbf{환경}(Environment): 에이전트가 상호작용하는 세계
\item \textbf{보상}(Reward): 좋은 행동 $\rightarrow$ +점수 / 나쁜 행동 $\rightarrow$ -점수
\end{itemize}

\vspace{0.3cm}
\emoji{point-right} 정답을 알려주지 않고, \textbf{스스로 최선의 행동}을 찾아가는 방식
\end{frame}

% ============================================================
% 슬라이드 21: 그림 08 — 강아지 훈련 비유 (columns: 그림+키워드)
% ============================================================
\begin{frame}{강화학습 비유: 강아지 훈련 \emoji{dog}}
\begin{columns}
\begin{column}{0.50\textwidth}
\includegraphics[width=\textwidth]{AI디지털리터러시_1주차2차시_19_강아지훈련_강화학습비유.png}

{\tiny 그림 08: 강아지훈련\_강화학습비유}
\end{column}
\begin{column}{0.45\textwidth}
\textbf{① 상황}: 강아지가 마당에 있음

\vspace{0.2cm}
\textbf{② 행동}: ``앉아!'' $\rightarrow$ 앉음

\vspace{0.2cm}
\textbf{③ 보상}: 잘했어! 간식! \emoji{bone}

\vspace{0.2cm}
\textbf{④ 학습}: 앉아 = 간식

\vspace{0.4cm}
\emoji{point-right} 다음번엔 바로 앉음!
\end{column}
\end{columns}
\end{frame}

% ============================================================
% 슬라이드 22: 그림 09 — 게임 점수 획득 (columns: 그림+키워드)
% ============================================================
\begin{frame}{강화학습 예시: 게임 AI \emoji{video-game}}
\begin{columns}
\begin{column}{0.40\textwidth}
\includegraphics[width=\textwidth]{AI디지털리터러시_1주차2차시_20_게임점수획득_강화학습과정.png}

{\tiny 그림 09: 게임점수획득\_강화학습과정}
\end{column}
\begin{column}{0.55\textwidth}
\textbf{100만 번 반복 학습}

\vspace{0.3cm}
\begin{tabular}{rl}
1,000번 & 50점 \emoji{loudly-crying-face} \\
10,000번 & 300점 \emoji{slightly-smiling-face} \\
100,000번 & 850점 \emoji{grinning-face} \\
1,000,000번 & 2,500점 \emoji{trophy} \\
\end{tabular}

\vspace{0.4cm}
\emoji{point-right} AI는 지치지 않고 24시간 연습!
\end{column}
\end{columns}
\end{frame}

% %%%%%%%%%%%%%%%%%%%%%%%%%%%%%%%%%%%%%%%%%%%%%%%%%%%%%%%%%%%%
% 전개 3: 머신러닝 4단계 프로세스
% %%%%%%%%%%%%%%%%%%%%%%%%%%%%%%%%%%%%%%%%%%%%%%%%%%%%%%%%%%%%
\section{전개 3: 머신러닝 4단계 프로세스}

% ============================================================
% 슬라이드 23: 그림 10 — 머신러닝 4단계 프로세스 (그림 단독)
% ============================================================
\begin{frame}{머신러닝 학습 과정 4단계}
\begin{center}
\includegraphics[width=0.85\textwidth]{AI디지털리터러시_1주차2차시_09_머신러닝4단계프로세스.png}

\vspace{0.1cm}
{\tiny 그림 10: 머신러닝\_4단계\_프로세스}
\end{center}
\end{frame}

% ============================================================
% 슬라이드 24: 그림 11 — 편향된 데이터 vs 좋은 데이터 (그림 단독)
% ============================================================
\begin{frame}{1단계 데이터 수집 — 데이터의 질이 중요!}
\begin{center}
\includegraphics[width=0.85\textwidth]{AI디지털리터러시_1주차2차시_13_편향된데이터vs좋은데이터.png}

\vspace{0.1cm}
{\tiny 그림 11: 편향된데이터vs좋은데이터}
\end{center}
\end{frame}

% ============================================================
% 슬라이드 25: 데이터 수집 3원칙 설명
% ============================================================
\begin{frame}{데이터 수집의 3가지 원칙}
\begin{center}
\begin{tikzpicture}
\node[draw=mainblue, fill=lightblue, rounded corners=8pt, inner sep=12pt, minimum width=3cm, text width=3cm, align=center] at (0,0) {\textbf{① 양}\\충분히 많은 데이터\\(수천\textasciitilde 수만 개)};
\node[draw=mainteal, fill=lightmint, rounded corners=8pt, inner sep=12pt, minimum width=3cm, text width=3cm, align=center] at (4.5,0) {\textbf{② 질}\\정확하고 깨끗한\\데이터};
\node[draw=mainpurple, fill=lightcream, rounded corners=8pt, inner sep=12pt, minimum width=3cm, text width=3cm, align=center] at (9,0) {\textbf{③ 다양성}\\다양한 상황의\\데이터};
\end{tikzpicture}
\end{center}

\vspace{0.5cm}
\emoji{warning} \textbf{편향된 데이터} $\rightarrow$ 편향된 AI!

\vspace{0.2cm}
\emoji{point-right} 실내 고양이만 학습하면 $\rightarrow$ 야외 고양이를 인식 못함
\end{frame}

% ============================================================
% 슬라이드 26: 그림 12 — 학습 곡선 (그림 단독)
% ============================================================
\begin{frame}{2단계 학습 — 에포크와 학습 곡선}
\begin{center}
\includegraphics[width=0.85\textwidth]{AI디지털리터러시_1주차2차시_12_학습곡선.png}

\vspace{0.1cm}
{\tiny 그림 12: 학습곡선 (에포크별 정확도 변화)}
\end{center}
\end{frame}

% ============================================================
% 슬라이드 27: 학습 핵심 용어 설명
% ============================================================
\begin{frame}{학습 핵심 용어}
\textbf{에포크 (Epoch)}
\begin{itemize}
\item 전체 데이터를 1회 학습하는 단위
\item 에포크 \emoji{up-arrow} $\Rightarrow$ 정확도 \emoji{up-arrow}
\end{itemize}

\vspace{0.4cm}
\textbf{가중치 (Weights)}
\begin{itemize}
\item AI가 학습하며 조정하는 수치
\item ``이 특징이 얼마나 중요한가?''를 숫자로 표현
\end{itemize}

\vspace{0.4cm}
\emoji{point-right} 3주차 엔트리 실습에서 에포크 변화를 직접 관찰!
\end{frame}

% ============================================================
% 슬라이드 28: 그림 13 — 데이터 분할 (그림 단독)
% ============================================================
\begin{frame}{3단계 평가 — 데이터 분할}
\begin{center}
\includegraphics[width=0.85\textwidth]{AI디지털리터러시_1주차2차시_11_데이터분할.png}

\vspace{0.1cm}
{\tiny 그림 13: 데이터분할 (Train 70\% / Validation 15\% / Test 15\%)}
\end{center}
\end{frame}

% ============================================================
% 슬라이드 29: 과적합 설명
% ============================================================
\begin{frame}{왜 데이터를 나눌까? — 과적합 방지}
\begin{columns}
\begin{column}{0.45\textwidth}
\begin{center}
\textbf{\color{mainorange}\emoji{cross-mark} 과적합 상태}

\vspace{0.3cm}
\begin{tikzpicture}
\node[draw=mainorange, fill=lightcream, rounded corners=5pt, inner sep=8pt, text width=4.5cm, align=center] {
훈련 데이터: 100점\\
새 데이터: 50점\\[0.2cm]
$\rightarrow$ 암기만, 이해 못함
};
\end{tikzpicture}
\end{center}
\end{column}
\begin{column}{0.45\textwidth}
\begin{center}
\textbf{\color{maingreen}\emoji{check-mark-button} 좋은 학습}

\vspace{0.3cm}
\begin{tikzpicture}
\node[draw=maingreen, fill=lightmint, rounded corners=5pt, inner sep=8pt, text width=4.5cm, align=center] {
훈련 데이터: 85점\\
새 데이터: 82점\\[0.2cm]
$\rightarrow$ 일반화 능력 우수
};
\end{tikzpicture}
\end{center}
\end{column}
\end{columns}

\vspace{0.4cm}
\begin{center}
\emoji{point-right} 시험 범위만 달달 외우면? $\rightarrow$ 응용 문제에서 틀림!
\end{center}
\end{frame}

% ============================================================
% 슬라이드 30: 4단계 예측
% ============================================================
\begin{frame}{4단계 예측 — 실전 적용}
\textbf{모델 배포}
\begin{itemize}
\item \textbf{①} 학습 완료된 모델 저장
\item \textbf{②} 앱/웹사이트에 적용
\item \textbf{③} 실시간 예측 시작
\end{itemize}

\vspace{0.4cm}
\textbf{지속적 개선}
\begin{itemize}
\item 새 데이터 수집 $\rightarrow$ 재학습 (Re-training)
\item 성능 모니터링 $\rightarrow$ 더 좋은 모델로 업데이트
\end{itemize}
\end{frame}

% %%%%%%%%%%%%%%%%%%%%%%%%%%%%%%%%%%%%%%%%%%%%%%%%%%%%%%%%%%%%
% 정리
% %%%%%%%%%%%%%%%%%%%%%%%%%%%%%%%%%%%%%%%%%%%%%%%%%%%%%%%%%%%%
\section{정리}

% ============================================================
% 슬라이드 31: 엔트리 AI 연결
% ============================================================
\begin{frame}{엔트리 AI 모델과의 연결}
\begin{center}
\textbf{3주차에 배울 엔트리 이미지 분류 모델 = 지도학습!}
\end{center}

\vspace{0.3cm}
\begin{center}
\begin{tikzpicture}
\node[draw=mainblue, fill=lightblue, rounded corners=5pt, inner sep=6pt, minimum width=2.2cm] (a) at (0,0) {\textbf{①} 사진 수집};
\node[draw=mainteal, fill=lightmint, rounded corners=5pt, inner sep=6pt, minimum width=2.2cm] (b) at (3,0) {\textbf{②} 라벨 붙이기};
\node[draw=mainpurple, fill=lightcream, rounded corners=5pt, inner sep=6pt, minimum width=2.2cm] (c) at (6,0) {\textbf{③} 학습 클릭};
\node[draw=maingreen, fill=lightmint, rounded corners=5pt, inner sep=6pt, minimum width=2.2cm] (d) at (9,0) {\textbf{④} 에포크 관찰};
\node[draw=mainorange, fill=lightcream, rounded corners=5pt, inner sep=6pt, minimum width=2.2cm] (e) at (12,0) {\textbf{⑤} 예측 테스트};
\draw[->,thick] (a) -- (b);
\draw[->,thick] (b) -- (c);
\draw[->,thick] (c) -- (d);
\draw[->,thick] (d) -- (e);
\end{tikzpicture}
\end{center}

\vspace{0.3cm}
\emoji{point-right} 오늘 배운 4단계 프로세스를 \textbf{직접 체험}하게 됩니다!
\end{frame}

% ============================================================
% 슬라이드 32: 오늘 배운 내용 정리
% ============================================================
\begin{frame}{오늘 배운 내용 정리}
\emoji{check-mark-button} \textbf{AI$\supset$ML$\supset$DL 포함 관계}
\begin{itemize}
\item 기계학습은 AI를 구현하는 방법
\end{itemize}

\vspace{0.3cm}
\emoji{check-mark-button} \textbf{3가지 학습 유형}
\begin{itemize}
\item 지도학습: 정답 있음 (분류/회귀)
\item 비지도학습: 정답 없음 (군집화)
\item 강화학습: 보상으로 학습 (시행착오)
\end{itemize}

\vspace{0.3cm}
\emoji{check-mark-button} \textbf{머신러닝 4단계}
\begin{itemize}
\item 데이터 수집 $\rightarrow$ 학습 $\rightarrow$ 평가 $\rightarrow$ 예측
\end{itemize}
\end{frame}

% ============================================================
% 슬라이드 33: 퀴즈 안내
% ============================================================
\begin{frame}{다음 주 퀴즈 안내 \emoji{memo}}
\begin{center}
\begin{tabular}{ll}
\textbf{시기} & 2주차 1차시 시작 (15분) \\
\textbf{범위} & 1주차 1차시 + 1주차 2차시 전체 \\
\textbf{형식} & 객관식 10문항 (5점) + 단답형 5문항 (5점) \\
\textbf{배점} & 총 10점 \\
\end{tabular}
\end{center}

\vspace{0.5cm}
\textbf{준비 방법}:
\begin{itemize}
\item 강의교재 복습
\item 핵심 개념 확인 (AI/ML/DL, 3가지 학습유형, 4단계 프로세스)
\end{itemize}
\end{frame}

% ============================================================
% 슬라이드 34: 마무리 (standout)
% ============================================================
\begin{frame}[standout]
\emoji{clapping-hands} 수고하셨습니다!

\vspace{0.5cm}
다음 주에 만나요 \emoji{waving-hand}
\end{frame}

\end{document}
